\documentclass{article}
\usepackage[utf8]{inputenc}
\usepackage{geometry}
\usepackage{enumitem}
\usepackage{hyperref}
\usepackage{xcolor}
\usepackage{graphicx}
\usepackage{array}
\usepackage{multirow}
\usepackage{amsmath, amssymb}
\usepackage{chemfig}
\usepackage{mhchem}
\usepackage{tikz}
\usepackage{setspace}
\usepackage{soul}
\usepackage{mathtools}
\usepackage{booktabs}
\usepackage{chemformula}
\usepackage{siunitx}
\usetikzlibrary{positioning, calc}
\usepackage{longtable}
\geometry{margin=1in}
\setlength{\parskip}{0.5em}
\usepackage{cancel}


\begin{document}

\begin{center}
    \vspace*{-1cm}
    \begin{tabular}{p{12cm} r}
        & \textbf{Name:} \underline{\hspace{5cm}} \\
    \end{tabular}

    \vspace{0.2cm}
    {\LARGE \textbf{Week 15.2 Worksheet}}
\end{center}    

\noindent\begin{tabular}{|p{0.9\linewidth}|}
\hline
\textbf{(1)} Which solution has the highest [OH$^-$]? \\[0.2cm]
\begin{minipage}{0.85\linewidth}
\begin{enumerate}[label=\alph*.]
\item pH = 6.10 \item pH = 8.23 \item pH = 9.50 \item pH = 3.33
\end{enumerate}
\end{minipage}\\[0.2cm]
\hline
\end{tabular}

\noindent\begin{tabular}{|p{0.9\linewidth}|}
\hline
\textbf{(2)} Which substance is a strong base? \\[0.2cm]
\begin{minipage}[t]{0.85\linewidth}
\begin{enumerate}[label=\alph*.]
\item NH$_3$ \item Ca(OH)$_2$ \item HF \item HNO$_2$
\end{enumerate}
\end{minipage}\\[0.2cm]
\hline
\end{tabular}

\noindent\begin{tabular}{|p{0.9\linewidth}|}
\hline
\textbf{(3)} Which of the following statements about salt hydrolysis is correct? \\[0.2cm]
\begin{minipage}{0.85\linewidth}
\begin{enumerate}[label=\alph*.]
\item Only salts of strong acids hydrolyze
\item Only salts of strong bases hydrolyze
\item Salts of weak acids or weak bases hydrolyze
\item No salts hydrolyze in water
\end{enumerate}
\end{minipage}\\[0.2cm]
\hline
\end{tabular}

\noindent\begin{tabular}{|p{0.9\linewidth}|}
\hline
\textbf{(4)} Which salt forms a basic solution upon hydrolysis? \\[0.2cm]
\begin{minipage}{0.85\linewidth}
\begin{enumerate}[label=\alph*.]
\item KClO \item NH$_4$Br \item NaNO$_3$ \item LiCl
\end{enumerate}
\end{minipage}\\[0.2cm]
\hline
\end{tabular}
\newpage

\noindent\begin{tabular}{|p{0.9\linewidth}|}
\hline
\textbf{(5)} Which describes a neutral solution? \\[0.2cm]
\begin{minipage}{0.85\linewidth}
\begin{enumerate}[label=\alph*.]
\item $[H_3O^+] > [OH^-]$
\item $[H_3O^+] < [OH^-]$
\item $[H_3O^+] = [OH^-]$
\item $[H_3O^+] = 0$
\end{enumerate}
\end{minipage}\\[0.2cm]
\hline
\end{tabular}

\noindent\begin{tabular}{|p{0.9\linewidth}|}
\hline
\textbf{(6)} A solution has [$H_3O^+$] = $3.20\times10^{-6}$ M. Its pH is closest to: \\[0.2cm]
\begin{minipage}{0.85\linewidth}
\begin{enumerate}[label=\alph*.]
\item 4.49 \item 5.49 \item 6.49 \item 7.49
\end{enumerate}
\end{minipage}\\[0.2cm]
\hline
\end{tabular}

\noindent\begin{tabular}{|p{0.9\linewidth}|}
\hline
\textbf{(7)} Which salt will produce a basic solution upon dissolution due to its parent species? \\[0.2cm]
\begin{minipage}{0.85\linewidth}
\begin{enumerate}[label=\alph*.]
\item NaF \item NH$_4$Br \item KCl \item Ca(NO$_3$)$_2$
\end{enumerate}
\end{minipage}\\[0.2cm]
\hline
\end{tabular}

\noindent\begin{tabular}{|p{0.9\linewidth}|}
\hline
\textbf{(8)} In the reaction HCOOH + H$_2$O $\rightleftharpoons$ HCOO$^-$ + H$_3$O$^+$, HCOOH is acting as a(n): \\[0.2cm]
\begin{minipage}{0.85\linewidth}
\begin{enumerate}[label=\alph*.]
\item Base \item Acid \item Salt \item Neutral compound
\end{enumerate}
\end{minipage}\\[0.2cm]
\hline
\end{tabular}
\newpage

\noindent\begin{tabular}{|p{0.9\linewidth}|}
\hline
\textbf{(9)} Which of the following solutions is most basic? \\[0.2cm]
\begin{minipage}{0.85\linewidth}
\begin{enumerate}[label=\alph*.]
\item pH = 9.00 \item pH = 11.00 \item pH = 7.00 \item pH = 5.00
\end{enumerate}
\end{minipage}\\[0.2cm]
\hline
\end{tabular}

\noindent\begin{tabular}{|p{0.9\linewidth}|}
\hline
\textbf{(10)} Which of the following best explains why pKa values are useful? \\[0.2cm]
\begin{minipage}{0.85\linewidth}
\begin{enumerate}[label=\alph*.]
\item They show how fast an acid reacts
\item They show the strength of an acid using a log scale
\item They show solubility of an acid
\item They show the molecular mass of an acid
\end{enumerate}
\end{minipage}\\[0.2cm]
\hline
\end{tabular}

\noindent\begin{tabular}{|p{0.9\linewidth}|}
\hline
\textbf{(11)} Which term describes an acid that ionizes partially in water? \\[0.2cm]
\begin{minipage}[t]{0.85\linewidth}
\begin{enumerate}[label=\alph*.]
\item Strong acid \item Weak acid \item Strong base \item Neutral compound
\end{enumerate}
\end{minipage}\\[0.2cm]
\hline
\end{tabular}

\noindent\begin{tabular}{|p{0.9\linewidth}|}
\hline
\textbf{(12)} Which description correctly identifies the parent acid of the salt sodium acetate (NaOAc)? \\[0.2cm]
\begin{minipage}{0.85\linewidth}
\begin{enumerate}[label=\alph*.]
\item HCl \item CH$_3$COOH \item HNO$_3$ \item HF
\end{enumerate}
\end{minipage}\\[0.2cm]
\hline
\end{tabular}
\newpage

\noindent\begin{tabular}{|p{0.9\linewidth}|}
\hline
\textbf{(13)} Which species is the conjugate acid of NH$_3$? \\[0.2cm]
\begin{minipage}{0.85\linewidth}
\begin{enumerate}[label=\alph*.]
\item NH$_2^-$ \item NH$_4^+$ \item N$^{3-}$ \item NH$_2$OH
\end{enumerate}
\end{minipage}\\[0.2cm]
\hline
\end{tabular}

\noindent\begin{tabular}{|p{0.9\linewidth}|}
\hline
\textbf{(14)} In the stepwise dissociation of H$_2$SO$_4$, which statement is correct? \\[0.2cm]
\begin{minipage}{0.85\linewidth}
\begin{enumerate}[label=\alph*.]
\item Both protons dissociate completely \item The first proton dissociates strongly, the second weakly \item Neither proton dissociates \item The second proton dissociates more strongly than the first
\end{enumerate}
\end{minipage}\\[0.2cm]
\hline
\end{tabular}

\noindent\begin{tabular}{|p{0.9\linewidth}|}
\hline
\textbf{(15)} For the equilibrium $\text{HA} + \text{H}_2\text{O} \rightleftharpoons \text{H}_3\text{O}^+ + \text{A}^-$, which expression represents $K_a$? \\[0.2cm]
\begin{minipage}[t]{0.85\linewidth}
\begin{enumerate}[label=\alph*.]
\item $\dfrac{[\text{HA}]}{[\text{H}_3\text{O}^+][\text{A}^-]}$ \item $\dfrac{[\text{H}_3\text{O}^+][\text{A}^-]}{[\text{HA}]}$ \item $[\text{H}_3\text{O}^+][\text{A}^-][\text{HA}]$ \item $\dfrac{1}{[\text{HA}]}$
\end{enumerate}
\end{minipage}\\[0.2cm]
\hline
\end{tabular}

\noindent\begin{tabular}{|p{0.9\linewidth}|}
\hline
\textbf{(16)} In the reaction $\text{NH}_3 + \text{H}_2\text{O} \rightarrow \text{NH}_4^+ + \text{OH}^-$, which species acts as the Brönsted-Lowry base?\\[0.2cm]
\begin{minipage}[t]{0.85\linewidth}
\begin{enumerate}[label=\alph*.]
\item NH$_3$
\item H$_2$O
\item NH$_4^+$
\item OH$^-$
\end{enumerate}
\end{minipage}\\[0.2cm]
\hline
\end{tabular}
\newpage

\noindent\begin{tabular}{|p{0.9\linewidth}|}
\hline
\textbf{(17)} Which species is amphiprotic? \\[0.2cm]
\begin{minipage}{0.85\linewidth}
\begin{enumerate}[label=\alph*.]
\item OH$^-$ \item HSO$_4^-$ \item Cl$^-$ \item NH$_4^+$
\end{enumerate}
\end{minipage}\\[0.2cm]
\hline
\end{tabular}

\noindent\begin{tabular}{|p{0.9\linewidth}|}
\hline
\textbf{(18)} A solution has pH = 5.00. What is the hydronium concentration? \\[0.2cm]
\begin{minipage}{0.85\linewidth}
\begin{enumerate}[label=\alph*.]
\item $1.00\times10^{-3}$ M \item $1.00\times10^{-5}$ M \item $1.00\times10^{-7}$ M \item $1.00\times10^{-9}$ M
\end{enumerate}
\end{minipage}\\[0.2cm]
\hline
\end{tabular}

\noindent\begin{tabular}{|p{0.9\linewidth}|}
\hline
\textbf{(19)} If a salt's anion is the conjugate base of a weak acid, what will be the pH of its aqueous solution? \\[0.2cm]
\begin{minipage}{0.85\linewidth}
\begin{enumerate}[label=\alph*.]
\item Acidic \item Basic \item Neutral \item Impossible to determine
\end{enumerate}
\end{minipage}\\[0.2cm]
\hline
\end{tabular}

\noindent\begin{tabular}{|p{0.9\linewidth}|}
\hline
\textbf{(20)} A weak acid has Ka = $3.00\times10^{-7}$. Which pKa is closest? \\[0.2cm]
\begin{minipage}{0.85\linewidth}
\begin{enumerate}[label=\alph*.]
\item 3.96 \item 4.22 \item 5.44 \item 6.80
\end{enumerate}
\end{minipage}\\[0.2cm]
\hline
\end{tabular}
\newpage

\noindent\begin{tabular}{|p{0.9\linewidth}|}
\hline
\textbf{(21)} A solution has a hydronium concentration of $1.0\times10^{-4}\ \text{M}$. What is its pH? \\[0.2cm]
\begin{minipage}[t]{0.85\linewidth}
\begin{enumerate}[label=\alph*.]
\item 2 \item 3 \item 4 \item 5
\end{enumerate}
\end{minipage}\\[0.2cm]
\hline
\end{tabular}

\noindent\begin{tabular}{|p{0.9\linewidth}|}
\hline
\textbf{(22)} Which statement correctly distinguishes dissociation from ionization? \\[0.2cm]
\begin{minipage}{0.85\linewidth}
\begin{enumerate}[label=\alph*.]
\item Dissociation involves breaking covalent bonds to form new ions, while ionization separates pre-existing ions in an ionic lattice.
\item Dissociation separates pre-existing ions from an ionic compound without altering their identity or overall charge, while ionization forms ions from a molecular substance through a chemical change.
\item Dissociation always produces H$_3$O$^+$, while ionization never does.
\item Dissociation only occurs in strong acids, while ionization only occurs in weak acids.
\end{enumerate}
\end{minipage}\\[0.2cm]
\hline
\end{tabular}

\noindent\begin{tabular}{|p{0.9\linewidth}|}
\hline
\textbf{(23)} A weak base has $[\text{OH}^-] = 1.00\times10^{-5}\ \text{M}$. What is the pOH? \\[0.2cm]
\begin{minipage}{0.85\linewidth}
\begin{enumerate}[label=\alph*.]
\item 3.00 \item 5.00 \item 7.00 \item 9.00
\end{enumerate}
\end{minipage}\\[0.2cm]
\hline
\end{tabular}

\noindent\begin{tabular}{|p{0.9\linewidth}|}
\hline
\textbf{(24)} For the equilibrium HF + H$_2$O $\rightleftharpoons$ H$_3$O$^+$ + F$^-$, which species is the conjugate base? \\[0.2cm]
\begin{minipage}{0.85\linewidth}
\begin{enumerate}[label=\alph*.]
\item HF \item H$_3$O$^+$ \item F$^-$ \item H$_2$O
\end{enumerate}
\end{minipage}\\[0.2cm]
\hline
\end{tabular}
\newpage

\noindent\begin{tabular}{|p{0.9\linewidth}|}
\hline
\textbf{(25)} A base that ionizes only slightly in water is classified as a: \\[0.2cm]
\begin{minipage}{0.85\linewidth}
\begin{enumerate}[label=\alph*.]
\item Strong base \item Weak base \item Neutral molecule \item Strong acid
\end{enumerate}
\end{minipage}\\[0.2cm]
\hline
\end{tabular}

\noindent\begin{tabular}{|p{0.9\linewidth}|}
\hline
\textbf{(26)} The parent acid of NO$_2^-$ is: \\[0.2cm]
\begin{minipage}{0.85\linewidth}
\begin{enumerate}[label=\alph*.]
\item HNO$_3$ \item HNO$_2$ \item HNO \item HNO$_4$
\end{enumerate}
\end{minipage}\\[0.2cm]
\hline
\end{tabular}

\noindent\begin{tabular}{|p{0.9\linewidth}|}
\hline
\textbf{(27)} A solution has pOH = 3.60. Determine the pH, [H$_3$O$^+$], and [OH$^-$] of the solution.  Ensure concentrations are in molarity and pH/pOH to two decimals. \\[4cm]
\hline
\end{tabular}

\noindent\begin{tabular}{|p{0.9\linewidth}|}
\hline
\textbf{(28)} Calcium hydroxide reacts with hydrochloric acid to form calcium chloride and water. Which statement best describes the acid-base behavior occurring in this reaction?\\[0.2cm]
Ca(OH)$_2$ + 2HCl $\rightarrow$ CaCl$_2$ + 2H$_2$O\\[0.2cm]
\begin{minipage}[t]{0.85\linewidth}
\begin{enumerate}[label=\alph*.]
\item HCl donates protons to OH$^-$ ions, forming water as Ca$^{2+}$ pairs with Cl$^-$
\item Ca(OH)$_2$ donates protons to HCl, forming hydrogen gas
\item Both reactants accept protons, preventing water formation
\item Ca$^{2+}$ acts as the proton donor while Cl$^-$ acts as the proton acceptor
\end{enumerate}
\end{minipage}\\[0.2cm]
\hline
\end{tabular}
\newpage

\noindent\begin{tabular}{|p{0.9\linewidth}|}
\hline
\textbf{(29)} Which statement correctly describes a strong base? \\[0.2cm]
\begin{minipage}{0.85\linewidth}
\begin{enumerate}[label=\alph*.]
\item Completely dissociates in water \item Partially dissociates in water \item Does not react with water \item Cannot accept protons
\end{enumerate}
\end{minipage}\\[0.2cm]
\hline
\end{tabular}

\noindent\begin{tabular}{|p{0.9\linewidth}|}
\hline
\textbf{(30)} A solution with pOH = 3.00 has what [OH$^-$]? \\[0.2cm]
\begin{minipage}{0.85\linewidth}
\begin{enumerate}[label=\alph*.]
\item $1.00\times10^{-3}$ M \item $1.00\times10^{-7}$ M \item $1.00\times10^{-11}$ M \item $1.00\times10^{-14}$ M
\end{enumerate}
\end{minipage}\\[0.2cm]
\hline
\end{tabular}

\noindent\begin{tabular}{|p{0.9\linewidth}|}
\hline
\textbf{(31)} Which statement best describes the logarithmic nature of pKa? \\[0.2cm]
\begin{minipage}{0.85\linewidth}
\begin{enumerate}[label=\alph*.]
\item A decrease of 1 pKa means the acid is ten times stronger
\item A decrease of 1 pKa means the acid is half as strong
\item A decrease of 1 pKa means no measurable change
\item A decrease of 1 pKa means the acid becomes more basic
\end{enumerate}
\end{minipage}\\[0.2cm]
\hline
\end{tabular}

\noindent\begin{tabular}{|p{0.9\linewidth}|}
\hline
\textbf{(32)} Which of the following is a triprotic acid? \\[0.2cm]
\begin{minipage}{0.85\linewidth}
\begin{enumerate}[label=\alph*.]
\item HCl \item H$_2$CO$_3$ \item H$_3$PO$_4$ \item HF
\end{enumerate}
\end{minipage}\\[0.2cm]
\hline
\end{tabular}
\newpage

\noindent\begin{tabular}{|p{0.9\linewidth}|}
\hline
\textbf{(33)} Which is the conjugate acid of CO$_3^{2-}$? \\[0.2cm]
\begin{minipage}{0.85\linewidth}
\begin{enumerate}[label=\alph*.]
\item CO$_2$ \item HCO$_3^-$ \item H$_2$CO$_3$ \item CO$_3$H$_2$
\end{enumerate}
\end{minipage}\\[0.2cm]
\hline
\end{tabular}

\noindent\begin{tabular}{|p{0.9\linewidth}|}
\hline
\textbf{(34)} Which salt will undergo acidic hydrolysis when dissolved in water? \\[0.2cm]
\begin{minipage}{0.85\linewidth}
\begin{enumerate}[label=\alph*.]
\item KNO$_3$ \item NH$_4$Cl \item NaClO$_3$ \item CaCl$_2$
\end{enumerate}
\end{minipage}\\[0.2cm]
\hline
\end{tabular}

\noindent\begin{tabular}{|p{0.9\linewidth}|}
\hline
\textbf{(35)} A solution has pH = 11.00 What is the hydroxide concentration? \\[0.2cm]
\begin{minipage}[t]{0.85\linewidth}
\begin{enumerate}[label=\alph*.]
\item $1.00\times10^{-3}$ M \item $1.00\times10^{-11}$ M \item $1.00\times10^{-7}$ M \item $1.00\times10^{-14}$ M
\end{enumerate}
\end{minipage}\\[0.2cm]
\hline
\end{tabular}

\noindent\begin{tabular}{|p{0.9\linewidth}|}
\hline
\textbf{(36)} A Brönsted-Lowry acid is defined as a substance that: \\[0.2cm]
\begin{minipage}[t]{0.85\linewidth}
\begin{enumerate}[label=\alph*.]
\item Donates a proton
\item Accepts a proton
\item Releases electrons
\item Accepts electrons
\end{enumerate}
\end{minipage}\\[0.2cm]
\hline
\end{tabular}
\newpage

\noindent\begin{tabular}{|p{0.9\linewidth}|}
\hline
\textbf{(37)} A weak acid has pKa = 4.25 and [HA] = 0.10 M. If [H$_3$O$^+$] = $5.6\times10^{-5}$ M, calculate pH, pOH, and [OH$^-$].  Ensure concentrations are in molarity and pH/pOH to two decimals. \\[4cm]
\hline
\end{tabular}

\noindent\begin{tabular}{|p{0.9\linewidth}|}
\hline
\textbf{(38)} In the reaction $\text{H}_2\text{PO}_4^-$ + H$_2$O \rightleftharpoons \text{HPO}_4^{2-} + \text{H}_3O^+$, the conjugate base is: \\[0.2cm]
\begin{minipage}[t]{0.85\linewidth}
\begin{enumerate}[label=\alph*.]
\item H$_2$PO$_4^-$
\item HPO$_4^{2-}$
\item H$_3$O$^+$
\item None of the above
\end{enumerate}
\end{minipage}\\[0.2cm]
\hline
\end{tabular}

\noindent\begin{tabular}{|p{0.9\linewidth}|}
\hline
\textbf{(39)} Which of the following best defines ionization for a weak acid in water? \\[0.2cm]
\begin{minipage}[t]{0.85\linewidth}
\begin{enumerate}[label=\alph*.]
\item Complete separation into ions \item Partial separation into ions \item Reaction with a base to form a salt \item Combination with water molecules without forming ions
\end{enumerate}
\end{minipage}\\[0.2cm]
\hline
\end{tabular}

\noindent\begin{tabular}{|p{0.9\linewidth}|}
\hline
\textbf{(40)} A weak acid has pKa = 3.97. Its solution has [H$_3$O$^+$] = $8.00\times10^{-5}$ M. Determine pH, pOH, and [OH$^-$].  Ensure concentrations are in molarity and pH/pOH to two decimals.\\[4cm]
\hline
\end{tabular}
\newpage

\noindent\begin{tabular}{|p{0.9\linewidth}|}
\hline
\textbf{(41)} Which pair represents a conjugate acid-base pair?\\[0.2cm]
\begin{minipage}[t]{0.85\linewidth}
\begin{enumerate}[label=\alph*.]
\item NH$_3$/NH$_4^+$
\item OH$^-$/O$^{2-}$
\item Na$^+$/Na
\item Cl$^-$/Cl$_2$
\end{enumerate}
\end{minipage}\\[0.2cm]
\hline
\end{tabular}

\noindent\begin{tabular}{|p{0.9\linewidth}|}
\hline
\textbf{(42)} Which describes the relationship between Ka and acid strength? \\[0.2cm]
\begin{minipage}{0.85\linewidth}
\begin{enumerate}[label=\alph*.]
\item Larger Ka means stronger acid \item Smaller Ka means stronger acid \item Ka does not relate to strength \item Ka only applies to bases
\end{enumerate}
\end{minipage}\\[0.2cm]
\hline
\end{tabular}

\noindent\begin{tabular}{|p{0.9\linewidth}|}
\hline
\textbf{(43)} If the pH of a solution increases from 8.00 to 10.00, the [H$^+$] changes by a factor of: \\[0.2cm]
\begin{minipage}{0.85\linewidth}
\begin{enumerate}[label=\alph*.]
\item 2x \item 10x \item 50x \item 100x
\end{enumerate}
\end{minipage}\\[0.2cm]
\hline
\end{tabular}

\noindent\begin{tabular}{|p{0.9\linewidth}|}
\hline
\textbf{(44)} A strong base solution has [OH$^-$] = $6.08\times10^{-4}$ M. Calculate pOH, pH, and [H$_3$O$^+$] to two decimal places, with concentrations in molarity. \\[4cm]
\hline
\end{tabular}
\newpage

\noindent\begin{tabular}{|p{0.9\linewidth}|}
\hline
\textbf{(45)} A solution with pH = 2.00 has a hydronium concentration that is how many times larger than a solution with pH = 4.00? \\[0.2cm]
\begin{minipage}{0.85\linewidth}
\begin{enumerate}[label=\alph*.]
\item 2x \item 10x \item 100x \item 1000x
\end{enumerate}
\end{minipage}\\[0.2cm]
\hline
\end{tabular}

\noindent\begin{tabular}{|p{0.9\linewidth}|}
\hline
\textbf{(46)} If an acid has $K_a = 1.00\times10^{-5}$, which of the following is closest to its pKa? \\[0.2cm]
\begin{minipage}{0.85\linewidth}
\begin{enumerate}[label=\alph*.]
\item 2.20 \item 3.50 \item 5.33 \item 7.25
\end{enumerate}
\end{minipage}\\[0.2cm]
\hline
\end{tabular}

\noindent\begin{tabular}{|p{0.9\linewidth}|}
\hline
\textbf{(47)} A solution has pH = 5.83. Find [H$_3$O$^+$], pOH, and [OH$^-$]. Concentrations in M, pH/pOH to two decimal places. \\[4cm]
\hline
\end{tabular}

\noindent\begin{tabular}{|p{0.9\linewidth}|}
\hline
\textbf{(48)} Which statement correctly defines ionization for a base in water? \\[0.2cm]
\begin{minipage}[t]{0.85\linewidth}
\begin{enumerate}[label=\alph*.]
\item Forming OH$^-$ by reacting chemically with water rather than releasing pre-existing ions
\item Separating into pre-existing ions without chemical change
\item Producing H$_3$O$^+$ by transferring a proton to water
\item Splitting into neutral molecules without forming ions
\end{enumerate}
\end{minipage}\\[0.2cm]
\hline
\end{tabular}
\newpage

\noindent\begin{tabular}{|p{0.9\linewidth}|}
\hline
\textbf{(49)} Which species will act as a base according to Brönsted-Lowry? \\[0.2cm]
\begin{minipage}{0.85\linewidth}
\begin{enumerate}[label=\alph*.]
\item H$_3$O$^+$ \item NH$_3$ \item HCl \item HNO$_3$
\end{enumerate}
\end{minipage}\\[0.2cm]
\hline
\end{tabular}

\noindent\begin{tabular}{|p{0.9\linewidth}|}
\hline
\textbf{(50)} A solution has [H$_3$O$^+$] = $2.22\times10^{-6}$ M. Calculate the pH, pOH, and [OH$^-$]. \\[4cm]
\hline
\end{tabular}

\noindent\begin{tabular}{|p{0.9\linewidth}|}
\hline
\textbf{(51)} Which of the following is a strong acid? \\[0.2cm]
\begin{minipage}[t]{0.85\linewidth}
\begin{enumerate}[label=\alph*.]
\item HF \item HCl \item HCN \item CH$_3$COOH
\end{enumerate}
\end{minipage}\\[0.2cm]
\hline
\end{tabular}

\noindent\begin{tabular}{|p{0.9\linewidth}|}
\hline
\textbf{(52)} A solution has [H$_3$O$^+$] = $4.58\times10^{-5}$ M. Calculate the pH, pOH, and [OH$^-$]. Express concentrations in molarity and pH/pOH to two decimal places. \\[4cm]
\hline
\end{tabular}
\newpage

\noindent\begin{tabular}{|p{0.9\linewidth}|}
\hline
\textbf{(53)} Which salt produces an acidic solution? \\[0.2cm]
\begin{minipage}{0.85\linewidth}
\begin{enumerate}[label=\alph*.]
\item NaCl \item KCN \item NH$_4$NO$_3$ \item NaOAc
\end{enumerate}
\end{minipage}\\[0.2cm]
\hline
\end{tabular}

\noindent\begin{tabular}{|p{0.9\linewidth}|}
\hline
\textbf{(54)} Which salt is formed from a strong acid and strong base and therefore produces a neutral solution? \\[0.2cm]
\begin{minipage}{0.85\linewidth}
\begin{enumerate}[label=\alph*.]
\item NaCl \item NH$_4$Cl \item NaOAc \item KCN
\end{enumerate}
\end{minipage}\\[0.2cm]
\hline
\end{tabular}

\noindent\begin{tabular}{|p{0.9\linewidth}|}
\hline
\textbf{(55)} When hydrochloric acid reacts with sodium hydroxide to form sodium chloride and water, which statement best describes the acid-base process taking place?\\[0.2cm]
HCl + NaOH $\rightarrow$ NaCl + H$_2$O\\[0.2cm]
\begin{minipage}[t]{0.85\linewidth}
\begin{enumerate}[label=\alph*.]
\item HCl donates a proton to OH$^-$, forming water as Na$^+$ pairs with Cl$^-$
\item NaOH donates protons to HCl, forming hydrogen gas
\item Both reactants act as proton donors, preventing neutralization
\item Cl$^-$ acts as the proton donor while Na$^+$ accepts the proton
\end{enumerate}
\end{minipage}\\[0.2cm]
\hline
\end{tabular}

\noindent\begin{tabular}{|p{0.9\linewidth}|}
\hline
\textbf{(56)} A weak acid HA has a pKa of 4.70. Which of the following pH values will favor the deprotonated form, A$^-$? \\[0.2cm]
\begin{minipage}{0.85\linewidth}
\begin{enumerate}[label=\alph*.]
\item 2.00 \item 3.00 \item 4.70 \item 7.00
\end{enumerate}
\end{minipage}\\[0.2cm]
\hline
\end{tabular}
\newpage

\noindent\begin{tabular}{|p{0.9\linewidth}|}
\hline
\textbf{(57)} According to the Arrhenius definition, which statement correctly describes an acid and a base?\\[0.2cm]
\begin{minipage}[t]{0.85\linewidth}
\begin{enumerate}[label=\alph*.]
\item An acid produces H$^+$ in water, and a base produces OH$^-$ in water
\item An acid donates a proton, and a base accepts a proton
\item An acid accepts an electron pair, and a base donates an electron pair
\item An acid and base both increase the concentration of ions equally
\end{enumerate}
\end{minipage}\\[0.2cm]
\hline
\end{tabular}

\noindent\begin{tabular}{|p{0.9\linewidth}|}
\hline
\textbf{(58)} In the equilibrium $\text{HF} + \text{H}_2\text{O} \rightleftharpoons \text{H}_3\text{O}^+ + \text{F}^-$, which is the conjugate acid of F$^-$? \\[0.2cm]
\begin{minipage}[t]{0.85\linewidth}
\begin{enumerate}[label=\alph*.]
\item HF
\item H$_3$O$^+$
\item H$_2$O
\item F$^-$ has no conjugate acid
\end{enumerate}
\end{minipage}\\[0.2cm]
\hline
\end{tabular}

\noindent\begin{tabular}{|p{0.9\linewidth}|}
\hline
\textbf{(59)} Which acid has the strongest conjugate base? \\[0.2cm]
\begin{minipage}{0.85\linewidth}
\begin{enumerate}[label=\alph*.]
\item HCl \item HNO$_3$ \item HF \item CH$_3$COOH
\end{enumerate}
\end{minipage}\\[0.2cm]
\hline
\end{tabular}

\noindent\begin{tabular}{|p{0.9\linewidth}|}
\hline
\textbf{(60)} Which statement is true about conjugate acid-base pairs? \\[0.2cm]
\begin{minipage}{0.85\linewidth}
\begin{enumerate}[label=\alph*.]
\item They differ by two protons \item They differ by one proton \item They differ by one electron \item They differ by two electrons
\end{enumerate}
\end{minipage}\\[0.2cm]
\hline
\end{tabular}
\end{document}



